\section{A Pareto-Improvement Counterfactual}\label{sec:pareto}

The 10--11 percent reduced-form price rise and the 29--45 percent welfare loss raise an obvious question: is there a cheaper rule that achieves the same SME-protection objective? Section~\ref{sec:welfare} supplied the answer arithmetically; this section develops it structurally.

A 10 percent price preference---of the kind the United States Small Business Administration has operated since the 1970s and that \citet{denes1997} documents in historical US federal procurement---works by scoring SME bids at a discount for purposes of winner selection while having the government pay the actual bid. Non-SMEs remain admissible and keep competing; SMEs get protection only to the extent that their bids fall within 10 percent of the non-SME minimum. Within the model's Bayes-Nash equilibrium, I simulate the preference rule on the same UH-clean cost distributions that anchor the set-aside decomposition. The result:
\begin{equation*}
\mathrm{V}_3 - \mathrm{V}_1 \;=\; +0.015 \text{ (non-pharma)}, \qquad -0.004 \text{ (pharma)},
\end{equation*}
versus $\mathrm{V}_0 - \mathrm{V}_1 = +0.245$ and $+0.308$ respectively under the full set-aside. The price cost to the government is essentially zero. The SME win-rate gain is of the same order as under the full set-aside---SMEs are selected as the preferred bidder in a comparable share of auctions---but the payment rule does not pull the equilibrium price up with them. Table~\ref{tab:v3_apv} reports the full comparison across four policy variants; Figure~\ref{fig:v3_welfare_weight} summarizes the welfare-weight implications.

\begin{figure}[!htbp]
\centering
\includegraphics[width=0.98\textwidth]{fig_v3_welfare_weight.pdf}
\caption[Implicit welfare weight to prefer the set-aside over the price preference.]{\textit{Implicit welfare weight $w^{\text{SME}}_\star$ required to prefer the full set-aside over a 10 percent price preference, by class and specification.} The grey band marks the 1.2--1.5 range typical of Brazilian cash-transfer programs; the dashed vertical line marks the utilitarian benchmark $w = 1$. In non-pharma, the implied weight sits above the transfer-program band under both the main equilibrium-selection specification ($w^{\text{SME}}_\star = 2.4$) and the strict-invariance benchmark ($w^{\text{SME}}_\star = 1.9$); the set-aside is welfare-dominated by the preference rule under any plausible distributional aggregator. In pharma, the implied weight under the main specification ($w^{\text{SME}}_\star = 3.0$) is far above the band, but reverses to 0.7 under strict invariance---below the utilitarian benchmark. The bifurcation is a property of the thin, heterogeneous pharma market, not a fragility of the result in non-pharma.}
\label{fig:v3_welfare_weight}
\end{figure}

\subsection{Welfare-weight identity}

The comparison formalizes through a standard welfare-weight identity. Let $w^{\text{SME}}$ denote the social weight on SME producer surplus. A planner is indifferent between V0 (full set-aside) and V3 (price preference) when
\begin{equation}\label{eq:welfare_weight}
w^{\text{SME}}_\star \;=\; \frac{\mathrm{Loss}(V_0) - \mathrm{Loss}(V_3)}{\mathrm{SMESurplus}(V_0) - \mathrm{SMESurplus}(V_3)}.
\end{equation}
With the point estimates of Section~\ref{sec:welfare} and the SME-surplus differentials implicit in V0--V3 ($\mathrm{Loss}(V_0) - \mathrm{Loss}(V_3) \approx 0.22$ in non-pharma and $0.30$ in pharma; $\mathrm{SMESurplus}(V_0) - \mathrm{SMESurplus}(V_3) \approx 0.09$ and $0.10$), the indifference weight is
\[
w^{\text{SME}}_\star \;\approx\; 2.4 \text{ (non-pharma)}, \qquad 3.0 \text{ (pharma)}.
\]
For the full set-aside to be welfare-optimal, the planner would need to value one real of SME producer surplus at 2.4 to 3.0 reais of general welfare. Brazilian cash-transfer programs (Bolsa Fam\'{\i}lia, Aux\'{\i}lio Brasil) have been defended under implicit weights in the 1.2--1.5 range---the shadow price of a real transferred to households in the lower tail of the income distribution, relative to a real in general tax revenue. The SMEs receiving set-aside benefits are medium-sized formal-sector firms, not households in the lower tail; the standard welfare argument for concentrating fiscal resources on them is \emph{weaker} than for cash transfers, not stronger. The implicit weight required to rationalize the current rule over the preference thus sits well above the weight Brazilian social policy places on population groups with the most established claim to such support. On these numbers the set-aside is welfare-inferior to the preference rule under any plausible distributional aggregator.

\subsection{Two implementation caveats}

The V3 counterfactual is Vickrey-equivalent. In a first-price sealed-bid auction with the same preference rule, equilibrium bidding would adjust in ways \citet{maskinriley2000} characterize: SMEs bidding more aggressively anticipating the preference, non-SMEs less aggressively anticipating the lost share. The direction of the adjustment typically reduces SME rents and increases the government's realized savings, so the Pareto-improvement claim carries through in sign even though the exact surplus split depends on the strategic response the paper does not solve. I flag this caveat explicitly; the preference-rule finding is directional rather than pinned to the Vickrey magnitude.

A preference rule is also procedurally more complex than a set-aside: the buyer must score bids twice and disclose the winner-selection rule. In the US SBA framework the rule has been operational for decades without significant friction. Brazilian procurement law already contemplates price preferences for several categories. The scale-up to Group-65 medical supplies would be a procedural rather than a conceptual move.

A back-of-envelope on the administrative cost of the preference rule sharpens the comparison. The US SBA HUBZone preference, which has run since 1997 with a 10 percent price preference on top of a set-aside, has been documented to add modest unit-administrative cost relative to a plain set-aside: scoring bids twice, preserving an audit trail, and adjudicating challenges from non-preferred bidders are real but bounded steps. Take an upper bound: assume a Brazilian preference rule costs three times as much per auction in administrative effort as the existing set-aside. With BEC processing roughly 800{,}000 procurement offers per year (Section~\ref{sec:institutional}) and a per-auction administrative cost of order R\$5--R\$15 under the existing set-aside, a tripling of that cost adds R\$10--R\$30 per auction in administrative loading. Against the 24 percent welfare cost differential the preference rule recovers in non-pharma at $\lambda = 0.30$---roughly R\$0.24 per real of procurement value---the administrative loading is two to three orders of magnitude smaller than the welfare gain, even at this conservative tripling assumption. The administrative cost is real but not load-bearing for the comparison.

\subsection{Is 10 percent the optimal preference rate?}\label{sec:optimal_preference}

The 10 percent rate is conventional---calibrated to the United States Small Business Administration HUBZone benchmark---but conventional is not optimal. To check, I simulate the welfare-and-win-rate frontier across preference rates from 0 to 30 percent in five-percentage-point increments. Table~\ref{tab:v3_preference_grid} reports the full grid; Figure~\ref{fig:v3_optimal_preference} plots the welfare-loss curve.

Two patterns emerge. First, the SME win-rate rises smoothly with the preference rate---from 17 percent at $k=0$ to 33 percent at $k=30$ in non-pharma; from 15 to 24 percent in pharma. The full set-aside delivers an SME win-rate of 100 percent, but the preference rule between 5 and 30 percent recovers a non-trivial share of the distributive gain. Second, the welfare cost of the preference rule rises only slowly until about 15--20 percent and then accelerates. Across the entire 0--30 grid, the welfare loss in non-pharma stays below 5 percent of $p_{S_1}$; in pharma it stays below 3 percent. The full set-aside, by contrast, delivers welfare losses of 28 percent in non-pharma and 44 percent in pharma. The order-of-magnitude gap between the preference rule and the set-aside is robust to the choice of $k$ within the 0--30 percent range.

The 10 percent rate is therefore not strictly optimal in the welfare-only sense. Rates closer to 5 percent in pharma deliver \emph{negative} welfare loss (the rule is welfare-improving) at modest SME win-rate gains, suggesting that the policy can in principle be designed to be Pareto-improving rather than merely Pareto-dominant over the set-aside. But across the policy-relevant range (5--15 percent), the choice of rate moves welfare cost by less than two percentage points of $p_{S_1}$. The substantive result is that the preference rule's welfare frontier sits an order of magnitude below the set-aside's at every $k$ in the policy-relevant range. The 10 percent rate is a defensible default, not a fragile choice.

\begin{figure}[!htbp]
\centering
\includegraphics[width=0.92\textwidth]{fig_v3_optimal_preference.pdf}
\caption[Welfare loss frontier under price-preference rates 0--30 percent.]{\textit{Welfare loss as \% of $p_{S_1}$, by class and price-preference rate $k \in [0, 30]$ percent.} Solid line plots simulated welfare loss under a price-preference of $k$ percent for SME bids. Open shapes mark the seven grid points. The horizontal dashed lines mark the welfare loss under the full set-aside (V0): 28 percent in non-pharma, 44 percent in pharma. The preference rule's welfare cost stays below 5 percent of $p_{S_1}$ in non-pharma and below 3 percent in pharma across the entire grid. Numbers from Table~\ref{tab:v3_preference_grid}.}
\label{fig:v3_optimal_preference}
\end{figure}

\subsection{Falsification test of Assumption~\ref{a:setaside}: strict primitive invariance}\label{sec:falsification}

Assumption~\ref{a:setaside} permits the post-period SME cost distribution $F_c^{\text{SME,Post}}$ to differ from the pre-period $F_c^{\text{SME,Pre}}$ as an equilibrium-selection object: the protected regime changes which SMEs find entry profitable. The headline welfare-dominance claim of this section depends on that assumption. A natural concern, anticipated by a careful reader, is that the assumption is unfalsifiable in the main specification---any inconvenient finding could be relabeled "selection". The paper offers an explicit falsification test: replace $F_c^{\text{SME,Post}}$ by $F_c^{\text{SME,Pre}}$ and re-run the welfare arithmetic. If the welfare-dominance claim survives, the result does not hinge on the selection assumption. If the claim reverses, the assumption is load-bearing.

Table~\ref{tab:v3_strict_invariance} reports the test. In non-pharma the preference rule continues to dominate the set-aside under strict invariance, with an implicit welfare weight of $w_\star^{\text{SME}} = 1.9$---above the Brazilian transfer-program benchmark and well above the utilitarian $w = 1$. The non-pharma welfare-dominance claim therefore passes the falsification test cleanly. In pharma the implicit weight falls to $w_\star^{\text{SME}} = 0.7$ and the ranking reverses: a utilitarian planner would prefer the set-aside to the preference rule under strict invariance.

I draw two conclusions from this. First, the welfare-dominance claim in non-pharma is robust across both specifications and is the paper's headline policy finding. The 10 percent price preference Pareto-dominates the SME-only set-aside in non-pharmaceutical procurement under both equilibrium-selection and strict-primitive-invariance specifications of Assumption~\ref{a:setaside}; this conclusion does not depend on the falsifiability of the selection mechanism. Second, the welfare-dominance claim in pharma is conditional on the equilibrium-selection assumption. Under strict invariance the conclusion reverses, and the paper does not claim otherwise. A reader who finds the strict-invariance benchmark more credible than the equilibrium-selection main specification should read the policy result as: \emph{the preference rule dominates in non-pharma; in pharma the comparison turns on whether the post-policy SME pool is the same primitives as the pre-policy pool (in which case the set-aside is preferred) or whether protection induces a different SME mix (in which case the preference rule is preferred)}. This is not a fragility of the methodology but a substantive feature of thin, heterogeneous markets, where the identity of the induced entrants is first-order for equilibrium price formation. The paper's contribution is to locate where the policy question becomes structural and to show that, where it does, it is the entrant-mix question specifically that drives the sensitivity.

Two further observations. The bid-distribution moments of Section~\ref{sec:rf} support the equilibrium-selection reading: the within-auction bid distribution shifts in pharma when the rule binds in ways that an unchanged $F_c^{\text{SME}}$ would not predict. The primitive-invariance KS test in Table~\ref{tab:v3_primitive_invariance} also rejects strict invariance precisely in the pharma SME stratum where one would expect entry selection to bite hardest, while it fails to reject in the strata where comparability is plausible. The data therefore weakly favor equilibrium selection over strict invariance in pharma; the strict-invariance benchmark serves as a stress test, not as the more credible specification. The headline non-pharma welfare-dominance claim does not require taking sides on this dispute.
